% !TeX root = ../../Book1.tex
% 纲要标题：### 21.3 经典方程
% 纲要位置：计划纲要.md，第 752 行。

\section{经典方程}
\label{b1:sec:2103}

二次、三次和四次方程的根式公式并非互不相关的技巧。它们背后的 Galois 群都嵌在可解的低阶对称群中；五次开始，对称群的可解性失效。不过，“一般方程没有通用根式公式”不表示每个具体五次方程都不能用根式求解。

% 纲要内容（仅作写作计划，尚非教材正文）：
%
% * 三次、四次方程的群论解释；从三次根的循环作用理解配对公式，从四根的三种配对构造辅助三次方程并联系 Ferrari 配方
% * 一般五次方程不可根式求解
% * 一个具体不可解五次方程：以 \(x^5-x-1\) 的模素数分解判定 Galois 群，引用第19.3节已证的循环型判据
% * 可解但非交换的 Galois 群：回引三次分裂域的群计算，解释正规子群链、单位根和根式塔
%

\subsection{三次、四次方程的群论解释}
\label{b1:subsec:2103:01}

设基域特征为零。若次数为 \(d\) 的多项式有 \(s\) 个不同的根，则其 Galois 群忠实作用在这 \(s\) 个根组成的集合上，因而嵌入 \(S_s\)，再将 \(S_s\) 看成 \(S_d\) 的子群即可。这里置换的是不同的根；即使按重数列出 \(d\) 个根，也不能据此增加域自同构。式\eqref{b1:eq:small-symmetric-derived}已经给出
\[
S_3\trianglerighteq A_3\trianglerighteq1,
\qquad S_4\trianglerighteq A_4\trianglerighteq V_4\trianglerighteq1
\]
都有交换因子。\cref{b1:thm:solvable-closure}把可解性传给子群，再由\cref{b1:thm:radical-solvability}得到三次、四次方程都根式可解。

\ProseParagraphBreak
可解性判据说明根式表达存在，却没有直接给出求根的步骤。三次方程的 Cardano 公式可以从根的置换中构造出来。将一般三次方程除以首项系数，再平移变量消去二次项，得到 \(y^3+py+q=0\)。设其三个根为 \(r_1,r_2,r_3\)，按重数列出，于是 \(r_1+r_2+r_3=0\)。先加入本原三次单位根 \(\omega\)。我们希望选取根的线性组合，使三轮换的作用成为乘以 \(\omega\) 或 \(\omega^2\)，从而把循环作用化为特征向量上的标量作用。下面两组系数正是为此选择的：
\[
u=\frac{r_1+\omega^2r_2+\omega r_3}{3},\qquad
v=\frac{r_1+\omega r_2+\omega^2r_3}{3}.
\]
由 \(1+\omega+\omega^2=0\)，反过来有
\[
r_1=u+v,\qquad r_2=\omega u+\omega^2v,\qquad
r_3=\omega^2u+\omega v.
\]
在固定 \(\omega\) 的条件下，循环改标 \((r_1,r_2,r_3)\mapsto(r_2,r_3,r_1)\) 使 \((u,v)\mapsto(\omega u,\omega^2v)\)，而交换 \(r_2,r_3\) 使 \(u,v\) 交换。因此 \(u^3,v^3\) 不受三轮换影响，却仍可能互换。这里的改标首先是对根的形式重排，并没有断言某个具体方程一定存在实现它的域自同构。对无重根方程，加入 \(\omega\) 后的 Galois 群中实际出现的置换才按这些公式作用；这个群可以严格小于 \(S_3\)。就上述形式置换而言，商群 \(S_3/A_3\cong C_2\) 记录 \(u^3,v^3\) 的交换，因而先用至多一次开平方将这两个量区分；随后 \(A_3\cong C_3\) 在 \(u,v\) 上分别乘以 \(\omega,\omega^2\)，对应于开立方恢复它们。

\ProseParagraphBreak
Vieta 公式给出 \(r_1+r_2+r_3=0\) 与 \(r_1r_2+r_1r_3+r_2r_3=p\)，所以
\[
r_1^2+r_2^2+r_3^2
=(r_1+r_2+r_3)^2-2(r_1r_2+r_1r_3+r_2r_3)=-2p.
\]
展开 \(u,v\) 的乘积并使用 \(1+\omega+\omega^2=0\)，便得
\[
9uv=r_1^2+r_2^2+r_3^2-(r_1r_2+r_1r_3+r_2r_3)
=-2p-p=-3p,
\]
所以 \(uv=-p/3\)，又有 \(u^3+v^3=r_1^3-3uvr_1=r_1^3+pr_1=-q\)。故 \(u^3,v^3\) 是二次方程
\[
T^2+qT-\frac{p^3}{27}=0
\]
的两个根。这就解释了 Cardano 公式中的平方根为何先于立方根出现。当 \(p=0\) 时直接开立方即可。当 \(p\neq0\) 时，设
\[
\Delta=\frac{q^2}{4}+\frac{p^3}{27},\qquad
U=-\frac q2+\sqrt\Delta,\qquad V=-\frac q2-\sqrt\Delta.
\]
有 \(U+V=-q\)、\(UV=-p^3/27\neq0\)。任取 \(u^3=U\)，再按乘积条件定义 \(v=-p/(3u)\)，则 \(v^3=V\)。于是
\[
(u+v)^3+p(u+v)+q=u^3+v^3+(3uv+p)(u+v)+q=0.
\]
利用已加入的 \(\omega\)，三个根为 \(\omega^j u+\omega^{-j}v\)、\(j=0,1,2\)，重根情形按重数理解。不能独立选取两个立方根后直接相加；\(uv=-p/3\) 是公式成立的必要配对条件。

\ProseParagraphBreak
实系数方程的计算途中也可能出现复数。例如 \(y^3-3y+1=0\) 的 \(\Delta=1/4-1=-3/4\)，选取复平方根后有 \(U=-1/2+i\sqrt3/2\)。因为 \(|U|=1\)，任一立方根 \(u\) 也满足 \(|u|=1\)，配对条件给出 \(v=1/u=\bar u\)。又因 \(\omega^{-j}=\overline{\omega^j}\)，三个表达式 \(\omega^j u+\omega^{-j}v\) 全为实数，计算却经过了复数中间量。这正说明根式可解允许使用复数中间量，并不要求每次开方都留在实域中。

\ProseParagraphBreak
四次方程也能通过一个辅助三次方程归结为二次方程。经平移后，它成为 \(y^4+py^2+qy+r=0\)。若 \(q=0\)，先将 \(y^2\) 作为未知量解二次方程。若 \(q\neq0\)，希望把四次式写成两个平方之差，从而分解为两个二次式。引入待定参数 \(m\)，有
\[
y^4+py^2+qy+r
=(y^2+m)^2-\bigl((2m-p)y^2-qy+m^2-r\bigr).
\]
要使括号中的二次式成为 \(\bigl(sy-q/(2s)\bigr)^2\)，比较二次项与常数项就要求
\[
s^2=2m-p,\qquad m^2-r=\frac{q^2}{4s^2}.
\]
消去 \(s\)，得到关于 \(m\) 的三次方程
\[
4(2m-p)(m^2-r)=q^2
\]
取它的任一根 \(m\)，因为 \(q\neq0\)，有 \(2m-p\neq0\)，可选取非零 \(s\) 满足 \(s^2=2m-p\)。三次方程保证上述常数项条件也成立，于是
\begin{align*}
y^4+py^2+qy+r
&=(y^2+m)^2-\Bigl(sy-\frac q{2s}\Bigr)^2\\
&=\Bigl(y^2-sy+m+\frac q{2s}\Bigr)
  \Bigl(y^2+sy+m-\frac q{2s}\Bigr).
\end{align*}
后面只剩两个二次方程。这说明四次求解为何能够归结为三次求解及连续开平方。

\ProseParagraphBreak
上述配方给出了两个二次因子，还可以从根的配对解释辅助三次方程的来源。设
\[
f(X)=X^4+bX^3+cX^2+dX+e
=\prod_{i=1}^4(X-\alpha_i),
\]
把四个根的指标分成两组，每组两个；不区分组内顺序，也不区分两组的先后。这三种无序划分为 \(12\mid34\)、\(13\mid24\)、\(14\mid23\)，分别对应下面三个配对量：
\[
\begin{aligned}
\beta_1&=\alpha_1\alpha_2+\alpha_3\alpha_4,\\
\beta_2&=\alpha_1\alpha_3+\alpha_2\alpha_4,\\
\beta_3&=\alpha_1\alpha_4+\alpha_2\alpha_3.
\end{aligned}
\]
若 \(e_i\) 表示四个根的第 \(i\) 个初等对称多项式，则 \(e_1=-b\)、\(e_2=c\)、\(e_3=-d\)、\(e_4=e\)。要找一个以三个 \(\beta_j\) 为根的基域多项式，需要把它们的初等对称式改写为 \(e_i\)。其中两两乘积之和展开后，恰含十二个形如 \(\alpha_i^2\alpha_j\alpha_k\)、\(i,j,k\) 互异的单项式，各出现一次；\(e_1e_3\) 包含同样的十二项，并另有四项 \(e_4\)，所以两者相差 \(4e_4\)。三重乘积的八项则可分成两类：
\[
\beta_1\beta_2\beta_3
=e_4\sum_{i=1}^4\alpha_i^2
 +\sum_{1\le i<j<k\le4}(\alpha_i\alpha_j\alpha_k)^2.
\]
第一个和为 \(e_1^2-2e_2\)；第二个和为 \(e_3^2-2e_2e_4\)，因为四个三重乘积的六个两两乘积之和恰为 \(e_2e_4\)。于是得到
\begin{align*}
\beta_1+\beta_2+\beta_3&=e_2=c,\\
\beta_1\beta_2+\beta_1\beta_3+\beta_2\beta_3
&=e_1e_3-4e_4=bd-4e,\\
\beta_1\beta_2\beta_3&=e_1^2e_4+e_3^2-4e_2e_4
=b^2e+d^2-4ce.
\end{align*}
因此它们满足同一个三次辅助方程，其多项式为
\begin{equation}\label{b1:eq:quartic-cubic-resolvent}
\begin{aligned}
R(T)&\coloneqq\prod_{j=1}^3(T-\beta_j)\\
&=T^3-cT^2+(bd-4e)T+4ce-b^2e-d^2.
\end{aligned}
\end{equation}
置换四个根就置换这三个量，这正是\subsecref{b1:subsec:0404:03}中三种二元划分的作用。它的核为
\[
V_4=\{1,(12)(34),(13)(24),(14)(23)\}.
\]
这些置换保持每一种无序划分的整体，允许交换组内的根，也允许交换两组。若四根互异，由
\[
\beta_1-\beta_2=(\alpha_1-\alpha_4)(\alpha_2-\alpha_3)
\]
及另两式可知三个 \(\beta_j\) 也互异。因此保持全部三种划分等价于固定全部三个配对量，\(V_4\) 在这里有了具体的域论含义。仍在四根互异的条件下，原四次方程的 Galois 群 \(G\leq S_4\) 在这些量上的作用核就是 \(G\cap V_4\)。

\ProseParagraphBreak
现在将配对量与配方中的参数对应起来。对前面的 \(y^4+py^2+qy+r\)，有 \(b=0\)、\(c=p\)、\(d=q\)、\(e=r\)，故
\[
R(2m)=4(2m-p)(m^2-r)-q^2.
\]
辅助方程使 \(R(2m)=0\)，而 \(R\) 的根正是三个 \(\beta_j\)，故 \(2m=\beta_j\) 对某个 \(j\) 成立。重新标号，把该配对写为 \(\{\alpha_1,\alpha_2\}\mid\{\alpha_3,\alpha_4\}\)，于是 \(2m=\alpha_1\alpha_2+\alpha_3\alpha_4\)。暂令 \(s_0=\alpha_1+\alpha_2=-(\alpha_3+\alpha_4)\)。四个交叉乘积之和为 \((\alpha_1+\alpha_2)(\alpha_3+\alpha_4)=-s_0^2\)，故 Vieta 公式给出
\[
p=\alpha_1\alpha_2+\alpha_3\alpha_4
  +(\alpha_1+\alpha_2)(\alpha_3+\alpha_4)=2m-s_0^2.
\]
再比较两组根组成的两个二次因子的一次项系数，得到
\[
s_0^2=2m-p,\qquad
q=s_0(\alpha_1\alpha_2-\alpha_3\alpha_4),\qquad
\alpha_1\alpha_2+\alpha_3\alpha_4=2m.
\]
因此 \(s_0=\pm s\)；若 \(s_0=-s\)，交换两组根后便有 \(s_0=s\)。当 \(q\neq0\) 时，\(s\neq0\)，两组根的积于是分别为 \(m+q/(2s)\) 与 \(m-q/(2s)\)，恰好得到上面两个二次因子。辅助三次方程确定配对量，开平方确定两组根的和，最后各解一个二次方程；群作用与求根计算在此描述的是同一过程。

\subsection{一般五次方程不可根式求解}
\label{b1:subsec:2103:02}

下面的“一般”有一个精确含义：把系数当作代数独立的变量，而不是先给它们某组具体数值。

\begin{theorem}[一般方程]\label{b1:thm:generic-equation}
设 \(n\geq1\) 为整数，并设 \(u_1,\ldots,u_n\) 在 \(\mathbb Q\) 上代数独立。多项式
\[
X^n-u_1X^{n-1}+u_2X^{n-2}-\cdots+(-1)^nu_n
\]
在 \(\mathbb Q(u_1,\ldots,u_n)\) 上的 Galois 群是 \(S_n\)。因此当 \(n\geq5\) 时，一般 \(n\) 次方程不根式可解。
\end{theorem}
\begin{proof}
先取代数独立的 \(t_1,\ldots,t_n\)，令 \(E=\mathbb Q(t_1,\ldots,t_n)\)，\(e_j\) 为这些变量的第 \(j\) 个初等对称多项式，并令 \(F=\mathbb Q(e_1,\ldots,e_n)\)。于是
\[
\prod_{i=1}^n(X-t_i)=X^n-e_1X^{n-1}+\cdots+(-1)^ne_n.
\]
每个 \(t_i\) 都是右端多项式的根，因而在 \(F\) 上代数。依次加入 \(t_1,t_2,\ldots,t_n\)，在已知一个根后除去对应线性因子，下一根满足的多项式次数就降低一。所以各步次数分别至多为 \(n,n-1,\ldots,1\)，由\cref{b1:thm:tower-law}得到 \([E:F]\leq n!\)，特别地，扩张有限。各 \(t_i\) 互异，故 \(E/F\) 是右端可分多项式的分裂域，为 Galois 扩张。

变量的任意置换保持多项式环的加法与乘法，并延拓为分式域 \(E\) 的自同构。每个 \(e_j\) 都是对称多项式，所以这些置换固定全部 \(e_j\)，也就逐点固定它们生成的系数域 \(F\)。不同置换在某个 \(t_i\) 上取值不同，给出 \(n!\) 个不同的 \(F\) 自同构。\cref{b1:thm:separable-embeddings}的嵌入数上界于是给出
\[
n!\leq|\operatorname{Aut}_F(E)|\leq[E:F]\leq n!.
\]
全部等号成立，所有自同构恰为这些变量置换，因此 Galois 群就是 \(S_n\)。

还需证明这些初等对称多项式代数独立。若 \(H\in\mathbb Q[U_1,\ldots,U_n]\) 满足 \(H(e_1,\ldots,e_n)=0\)，因 \(t_1,\ldots,t_n\) 代数独立，这就是关于 \(t_i\) 的多项式恒等式。\cref{b1:thm:fundamental-symmetric-polynomials}的唯一性具体说：一个对称多项式写成初等对称多项式的多项式表达时，所用的表达式唯一。按 \(t_1>\cdots>t_n\) 的词典序，\(e_1^{a_1}\cdots e_n^{a_n}\) 的首项指数依次为 \(a_1+\cdots+a_n,a_2+\cdots+a_n,\ldots,a_n\)，依次作差便恢复各 \(a_i\)，所以不同单项式的首项也不同。现在零多项式同时由 \(H\) 与零表达，所以 \(H=0\)，证明了 \(e_1,\ldots,e_n\) 的代数独立性。由 \(u_i\mapsto e_i\) 得到系数域 \(\mathbb Q(u_1,\ldots,u_n)\cong F\) 的同构，并把题述多项式送到上面已经考察的多项式，因此两者的 Galois 群同构。

当 \(n\geq5\) 时，\cref{b1:cor:symmetric-nonsolvable}说明 \(S_n\) 不可解；\cref{b1:thm:radical-solvability}遂排除根式求解。
\end{proof}

这就是\term[Abel-Ruffini-dingli]{Abel--Ruffini 定理}[Abel--Ruffini theorem]的通用方程形式。这里的“一般”始终指系数域 \(\mathbb Q(u_1,\ldots,u_n)\) 上那一个以代数独立变量为系数的多项式，不是关于随机选取数值系数的概率陈述。它排除的是以这些独立系数为起点的通用根式表达，并不排除特殊方程。例如 \(X^5-1\) 的分裂域是分圆域，Galois 群 \((\mathbb Z/5\mathbb Z)^\times\cong C_4\) 可解，因而它的全部根都能用根式表示。


\subsection{一个不可根式求解的五次方程}
\label{b1:subsec:2103:04}

一般方程的结论不能直接判定某一个给定方程。现在考察序言中的 \(X^5-X-1\)，用\cref{b1:lem:reduction-cycle-type}把模素数分解得到的循环型与群的传递性结合起来，确定它的 Galois 群。


\begin{proposition}\label{b1:prop:quintic-x5-x-1}
多项式 \(f=X^5-X-1\) 在 \(\mathbb Q\) 上的 Galois 群为 \(S_5\)，因而不根式可解。
\end{proposition}
\begin{proof}
先证明不可约性。五次多项式若可约，就至少有两个不可约因子；若每个因子的次数都不小于三，总次数至少为六，矛盾。因此只须排除一次与不可约二次因子。模三时，\(f\) 在 \(0,1,2\) 处的值都为二，没有一次因子。\(\mathbb F_3\) 上的首一不可约二次多项式只有
\(X^2+1\)、\(X^2+X+2\)、\(X^2+2X+2\)：对常数项非零的六个首一二次式逐一检查是否在 \(0,1,2\) 处有根，即得此表。用它们分别除 \(\bar f\)，余式依次为
\[
2,\qquad X+2,\qquad X+2.
\]
所以也没有不可约二次因子，\(\bar f\) 不可约；\cref{b1:prop:reduction-irreducible}说明 \(f\) 在 \(\mathbb Q\) 上不可约。

模二时，直接相乘可核验
\[
\bar f=(X^2+X+1)(X^3+X^2+1).
\]
两个因子在 \(\mathbb F_2\) 中均无根，次数分别为二和三，因而均不可约且互异，所以约化无重因子。有限域上的不可约多项式都可分，故约化也无重根，满足\cref{b1:lem:reduction-cycle-type}的条件。该引理给出群中的一个 \((2)(3)\) 型置换 \(\tau=(a\,b)(d\,e\,h)\)。这两个循环不相交，彼此交换，故
\[
\tau^3=(a\,b)^3(d\,e\,h)^3=(a\,b),
\]
得到群中的一个换位。

在分裂域中，若 \(a,b\) 是同一不可约因子的两个根，\cref{b1:lem:simple-embedding-extension}对基域恒等映射给出 \(K(a)\to K(b)\)、\(a\mapsto b\) 的同构；\cref{b1:thm:embedding-extension}延拓这个同构，正规性使得到的嵌入成为分裂域的自同构。因此一个不可约因子的全部根恰组成一个轨道。现在 \(f\) 不可约，故五个根只有一个轨道，群的作用传递。再由\cref{b1:thm:orbit-stabilizer}，五整除群阶。\cref{b1:thm:cauchy}给出五阶元素，在五个点上它必为五轮换 \(c\)。

将五个点按 \(c\) 标为 \(\mathbb Z/5\mathbb Z\)，使 \(c\) 的作用为加一。设刚找到的换位交换 \(i,j\)，记非零差为 \(d=j-i\)。用 \(c\) 的各次幂共轭这个换位，就得到所有 \((k\,\,k+d)\)。由于五为素数，从任一点反复加 \(d\)，依次经过全部五个点；所以任意两点都可用这些相邻关系连成一条不重复顶点的路径。设路径为 \(v_0,\ldots,v_r\)，它的各相邻换位都在群中。由第一个换位开始，若已经得到 \((v_0\,v_j)\)，就有
\[
(v_0\,v_j)(v_j\,v_{j+1})(v_0\,v_j)=(v_0\,v_{j+1})
\qquad(1\le j<r).
\]
沿路径归纳，得到两端的换位 \((v_0\,v_r)\)。任意两点均可作两端，故群包含全部换位；换位生成 \(S_5\)，于是原 Galois 群等于 \(S_5\)。由\cref{b1:cor:symmetric-nonsolvable}及\cref{b1:thm:radical-solvability}，方程不根式可解。
\end{proof}
这个具体判定解释了序言中 \(X^5-1\) 与 \(X^5-X-1\) 的差别：前者有交换的分圆 Galois 群，后者的群为不可解的 \(S_5\)。次数同为五，并不意味着两者有相同的根式可解性。

\ProseParagraphBreak
对可约多项式，能求出一个根也不意味着全部根可用根式求出。例如 \((X-1)(X^5-X-1)\) 有有理根 \(1\)，但它的五次因子已由上面的命题证明不根式可解。

\subsection{可解但非交换的 Galois 群}
\label{b1:subsec:2103:03}

回到 \(X^3-2\)。令 \(a=\sqrt[3]{2}\)、\(\zeta=\zeta_3\)、\(L=\mathbb Q(a,\zeta)\)。\cref{b1:prop:cubic-splitting-field-correspondence}已确定 \([L:\mathbb Q]=6\)、\(\operatorname{Gal}(L/\mathbb Q)\cong S_3\) 及其全部中间域。这里利用这些结论，解释根式求解怎样反映在子群与域塔上。

\ProseParagraphBreak
由 Galois 对应，正规子群链和相应的域塔为
\[
S_3\trianglerighteq A_3\trianglerighteq1,
\qquad
\mathbb Q\subseteq\mathbb Q(\zeta)\subseteq L.
\]
两步扩张的 Galois 群分别是 \(S_3/A_3\cong C_2\) 与 \(A_3\cong C_3\)。第一步通过
\[
\zeta=\frac{-1+\sqrt{-3}}2,
\qquad \mathbb Q(\zeta)=\mathbb Q(\sqrt{-3})
\]
开平方完成，反向关系 \(\sqrt{-3}=2\zeta+1\) 也说明两域相等。这一步使用复平方根，已离开实域，仍属于允许的根式运算。第二步加入 \(a\)，其立方 \(a^3=2\) 已在基域中。于是这条对应于正规子群链的域塔本身就是一条根式塔。

\ProseParagraphBreak
先加入单位根，使第二步也能由 Galois 理论描述：在 \(\mathbb Q(\zeta)\) 上，加入 \(a\) 便同时包含 \(a,\zeta a,\zeta^2a\) 三个共轭根，循环自同构把 \(a\) 送到 \(\zeta a\)。若先走 \(\mathbb Q\subseteq\mathbb Q(a)\)，虽然已经通过开立方取得一个根，这一步却不正规，因为另外两个根不在实域 \(\mathbb Q(a)\) 中。

\ProseParagraphBreak
若 \(\rho(a)=\zeta a\)、\(\rho(\zeta)=\zeta\)，而 \(\kappa\) 为复共轭，则 \(\kappa\rho(a)=\zeta^2a\)、\(\rho\kappa(a)=\zeta a\)，两者不同，故全群 \(S_3\) 非交换。但上述两层商群 \(S_3/A_3\cong C_2\) 与 \(A_3\cong C_3\) 都是交换群，所以 \(S_3\) 仍然可解。交换群本身总可解，而这个例子说明可解群可以非交换；根式求解所需要的是逐层出现交换的商群，不要求整个 Galois 群交换。
